\clearpage
\begin{explanationbox}
\textbf{先看具体比值：固定 $A(-1,0),B(1,0)$，要求 $PA=2PB$。}
\begin{center}\begin{tikzpicture}[font=\small]
\draw[black!30,rounded corners=2pt] (0,0) rectangle (14.4,8.2);
\node[anchor=west,font=\small\bfseries] at (.5,7.72) {定比为 2：圆周上的点都满足 $PA=2PB$};
\begin{scope}[shift={(5,3.7)},scale=1.55]
\coordinate (A) at (-1.0000000000,0.0000000000);
\coordinate (B) at (1.0000000000,0.0000000000);
\coordinate (C) at (1.6666666667,0.0000000000);
\coordinate (D) at (0.3333333333,0.0000000000);
\coordinate (E) at (3.0000000000,0.0000000000);
\coordinate (P) at (1.6666666667,1.3333333333);
\draw[blue,very thick] (C) circle (1.3333333333);
\draw[black!55] (-1.5,0)--(3.4,0);
\draw[blue,thick] (A)--(P);\draw[teal!80!black,thick] (B)--(P);\draw[orange!85!black,dashed] (C)--(P);
\fill[black] (A) circle (1.5pt) node[below=6pt] {$A$};
\fill[black] (B) circle (1.5pt) node[below=6pt] {$B$};
\fill[black] (C) circle (1.5pt) node[below=16pt] {$C$};
\fill[black] (P) circle (1.5pt) node[above=6pt] {$P$};
\coordinate (Q) at (1.2106398089,-1.2529234944);
\fill[black] (Q) circle (1.5pt) node[below=5pt] {$Q$};

\end{scope}
\node[anchor=west] at (.8,6.4) {$A(-1,0),\ B(1,0)$};\node[anchor=west] at (.8,5.6) {$PA=2PB$};\node[anchor=west] at (.8,1) {$C(5/3,0),\quad R=4/3$};
\end{tikzpicture}\end{center}

图中 $P=(5/3,4/3)$，有 $PA^2=80/9,PB^2=20/9$，所以比值确为 $2$。
$Q$ 代表同一圆周上的另一个点；换位置不会改变距离比。
距离更短的一端是 $B$，所以整圆位于 $AB$ 中垂线的 $B$ 侧。
这不是“以 $A$ 或 $B$ 为圆心画圆”：两个定点均不在轨迹圆周上。

\tcbbreak
\textbf{找圆不一定先背公式：先找直线 $AB$ 上的两个特殊点。}
\begin{center}\begin{tikzpicture}[font=\small]
\draw[black!30,rounded corners=2pt] (0,0) rectangle (14.4,8.2);
\node[anchor=west,font=\small\bfseries] at (.5,7.72) {同一距离比的内分点 $D$ 与外分点 $E$，恰是直径两端};
\begin{scope}[shift={(5,3.5)},scale=1.65]
\coordinate (A) at (-1.0000000000,0.0000000000);
\coordinate (B) at (1.0000000000,0.0000000000);
\coordinate (C) at (1.6666666667,0.0000000000);
\coordinate (D) at (0.3333333333,0.0000000000);
\coordinate (E) at (3.0000000000,0.0000000000);
\coordinate (P) at (1.6666666667,1.3333333333);
\draw[blue,very thick] (C) circle (1.3333333333);
\draw[black!55] (-1.5,0)--(3.4,0);
\draw[orange!85!black,very thick] (D)--(E);\fill[black] (A) circle (1.5pt) node[below=6pt] {$A$};
\fill[black] (B) circle (1.5pt) node[above=6pt] {$B$};
\fill[black] (C) circle (1.5pt) node[below=17pt] {$C$};
\fill[black] (D) circle (1.5pt) node[below left=6pt] {$D$};
\fill[black] (E) circle (1.5pt) node[below right=6pt] {$E$};

\end{scope}
\node[anchor=west] at (.8,6.5) {$AD/DB=AE/EB=2$};\node[anchor=west] at (.8,1.1) {$D=(1/3,0),\quad E=(3,0)$};\node[anchor=west] at (.8,.5) {$C$ 是 $DE$ 中点，$R=DE/2$};
\end{tikzpicture}\end{center}

内分点 $D$ 在 $A,B$ 之间，外分点 $E$ 在线段外；本图两点的距离比都为 $2$。
一旦知道它们是直径端点，圆心取中点、半径取端点距离的一半即可。
图中的 $B$ 在圆内，但\textbf{内部的点不因此成为轨迹点}。

\tcbbreak
\textbf{固定定点和单位，只把比值 $2$ 换成 $1/2$。}
\begin{center}\begin{tikzpicture}[font=\small]
\draw[black!30,rounded corners=2pt] (0,0) rectangle (14.4,8.2);
\node[anchor=west,font=\small\bfseries] at (.5,7.72) {固定 $A(-1,0),B(1,0)$：比值取倒数，圆关于中垂线对称};
\begin{scope}[shift={(7.2,3.6)},scale=1.6]
\coordinate (A) at (-1.0000000000,0.0000000000);
\coordinate (B) at (1.0000000000,0.0000000000);
\coordinate (C) at (1.6666666667,0.0000000000);
\coordinate (D) at (0.3333333333,0.0000000000);
\coordinate (E) at (3.0000000000,0.0000000000);
\coordinate (P) at (1.6666666667,1.3333333333);
\draw[-Stealth,black!50] (-3.6,0)--(3.6,0);\draw[blue,very thick] (C) circle (1.3333333333);
\draw[purple,thick,dashed] (-1.6666666667,0) circle (1.3333333333);\draw[black!40,dotted] (0,-1.7)--(0,1.7);\fill[black] (A) circle (1.5pt) node[below=5pt] {$A$};
\fill[black] (B) circle (1.5pt) node[below=5pt] {$B$};

\end{scope}
\node[text=purple] at (3.6,6.6) {$PA/PB=1/2$};\node[text=blue] at (10.8,6.6) {$PA/PB=2$};\node at (3.6,1) {$C=(-5/3,0)$};\node at (10.8,1) {$C=(5/3,0)$};
\end{tikzpicture}\end{center}

两个圆半径相同、位置互为镜像。$PA/PB<1$ 表示靠近 $A$，$PA/PB>1$ 表示靠近 $B$。
因此\textbf{先说清比值的顺序}，再判断圆在中垂线哪一侧。

\tcbbreak
\textbf{分界值单独看：$\lambda=1$ 是等距条件。}
\begin{center}\begin{tikzpicture}[font=\small]
\draw[black!30,rounded corners=2pt] (0,0) rectangle (14.4,8.2);
\node[anchor=west,font=\small\bfseries] at (.5,7.72) {$\lambda=1$：等距点组成 $AB$ 的中垂线};
\begin{scope}[shift={(7.2,3.6)},scale=1.6]
\coordinate (A) at (-1.0000000000,0.0000000000);
\coordinate (B) at (1.0000000000,0.0000000000);
\coordinate (P) at (0.0000000000,1.4000000000);
\draw[black!50] (-2.8,0)--(2.8,0);\draw[orange!85!black,very thick] (0,-1.45)--(0,1.9);\draw[blue,thick] (A)--(P);\draw[teal!80!black,thick] (B)--(P);\fill[black] (A) circle (1.5pt) node[below=6pt] {$A$};
\fill[black] (B) circle (1.5pt) node[below=6pt] {$B$};
\fill[black] (P) circle (1.5pt) node[above right=6pt] {$P$};

\end{scope}
\node[anchor=west] at (.8,6.5) {$PA=PB$};\node at (7.2,.6) {轨迹为 $x=0$，不是一个有限半径的圆};
\end{tikzpicture}\end{center}

标准坐标下，$PA^2-PB^2=4ax$。当距离相等时只能有 $x=0$；
反过来 $x=0$ 上每一点到 $A,B$ 都等距。
因此应直接得到中垂线，不能把 $\lambda=1$ 代入含 $\lambda^2-1$ 分母的圆公式。
\end{explanationbox}
